\documentclass[12pt, a4paper]{article}

\usepackage[T1]{fontenc}
\usepackage[utf8]{inputenc}
\usepackage{lmodern}
\usepackage[top=2.5cm, bottom=2.5cm, left=3cm, right=2.5cm]{geometry}
\usepackage{setspace}
\onehalfspacing

\usepackage{amsmath}
\usepackage{amssymb}
\usepackage{booktabs}
\usepackage{caption}
\usepackage{float}
\usepackage{xcolor}
\usepackage{listings}
\usepackage{microtype}
\usepackage{parskip}
\usepackage[colorlinks=true, linkcolor=black, urlcolor=blue]{hyperref}

% code style
\lstset{
  basicstyle=\ttfamily\small,
  backgroundcolor=\color{gray!8},
  frame=single,
  framerule=0pt,
  xleftmargin=6pt,
  xrightmargin=6pt,
  breaklines=true,
  columns=fullflexible,
}

\title{
  {\large Working Notes} \\[0.4cm]
  {\LARGE\bfseries LSTM and MIDAS: How the Notebooks Work}\\[0.4cm]
  {\normalsize Thesis — Forecasting Silver Prices}
}
\author{Asier Ugartechea Perez}
\date{\today}

\begin{document}
\maketitle
\tableofcontents
\newpage

% ─────────────────────────────────────────────────────────────────────────────
\section{Purpose of this document}
% ─────────────────────────────────────────────────────────────────────────────

This document walks through the logic of two notebook families that differ from the
ARIMA/VAR/Random Forest/XGBoost notebooks in their structure and implementation.
The ARIMA and tree-based notebooks follow a nearly identical template; LSTM and MIDAS
each have design decisions that are worth understanding before modifying them or
drawing conclusions from their output.

% ─────────────────────────────────────────────────────────────────────────────
\section{LSTM (\texttt{weekly/06\_lstm.ipynb})}
% ─────────────────────────────────────────────────────────────────────────────

\subsection{What it does at a high level}

The LSTM notebook trains a neural network that looks at the last 20 trading days of
features---prices, market returns, and optionally sentiment---and predicts the next
day's silver log-return. It tests four variants of increasing richness, mirroring the
ARIMA ablation:

\begin{table}[H]
\centering
\caption{LSTM variants and their ARIMA equivalents}
\begin{tabular}{llp{5.5cm}}
\toprule
Variant & Features & ARIMA analogue \\
\midrule
\texttt{LSTM-Y}      & Silver return only       & ARIMA (univariate) \\
\texttt{LSTM-EXOG}   & Silver + market variables & ARIMAX (no sentiment) \\
\texttt{LSTM-REDDIT} & EXOG + Reddit sentiment  & ARIMAX + Reddit \\
\texttt{LSTM-NEWS}   & EXOG + News sentiment    & ARIMAX + News \\
\bottomrule
\end{tabular}
\end{table}

\subsection{Architecture}

The model is a single-layer LSTM followed by a linear output layer:

\begin{equation}
  \hat{r}_{t+1} = W_{\text{out}} \cdot h_t^{(L)} + b
\end{equation}

\noindent where $h_t^{(L)}$ is the hidden state of the last LSTM layer at the final
timestep of the input sequence. An LSTM maintains a hidden state and a cell state that
it updates at each timestep, allowing information from earlier in the sequence to
influence the prediction---unlike a feedforward network, which treats each timestep
independently.

In code:

\begin{lstlisting}[language=Python]
class LSTMForecaster(nn.Module):
    def __init__(self, input_size, hidden_size=64, num_layers=2, dropout=0.2):
        super().__init__()
        self.lstm = nn.LSTM(input_size, hidden_size, num_layers=num_layers,
                            batch_first=True, dropout=dropout)
        self.fc = nn.Linear(hidden_size, 1)

    def forward(self, x):
        out, _ = self.lstm(x)
        return self.fc(out[:, -1, :])  # use only the last timestep's hidden state
\end{lstlisting}

\noindent \texttt{out[:, -1, :]} extracts the hidden state at the last position of
the sequence (day $t$), which the linear layer maps to a scalar return forecast.

\subsection{Input sequences}

The function \texttt{make\_sequences} slides a window of length
\texttt{SEQ\_LEN = 20} across the data:

\begin{itemize}
  \item Input: rows $t - 20$ through $t - 1$ (20 days of history, all past)
  \item Target: value at row $t$ (the next day's silver return)
\end{itemize}

This is strictly causal---no look-ahead bias. With \texttt{HORIZON = 1}, the target
is simply $r_t$ and the inverse-scaling back to log-return units is:

\begin{equation}
  \hat{r}_t = \hat{z}_t \cdot \sigma + \mu
\end{equation}

\noindent where $\hat{z}_t$ is the network's output in standardised units and
$(\mu, \sigma)$ are the training-set mean and standard deviation of the silver return,
estimated with \texttt{StandardScaler}.

\subsection{Training loop}

Each variant is trained independently (fresh weights, fresh scaler):

\begin{enumerate}
  \item \textbf{Early stopping}: training halts if validation loss does not improve
    for \texttt{PATIENCE = 10} consecutive epochs. The best checkpoint is saved and
    reloaded for inference---so the reported metrics reflect the best validation-loss
    weights, not the last epoch.
  \item \textbf{Learning rate scheduler}: \texttt{ReduceLROnPlateau} halves the
    learning rate when validation loss plateaus for 3 epochs, helping escape flat
    regions of the loss landscape.
  \item \textbf{Gradient clipping}: gradients are clipped at norm 1.0 to prevent
    the exploding-gradient problem that can destabilise LSTM training.
\end{enumerate}

\subsection{Why LSTM underperforms ARIMAX}

This is the most important thing to understand about the results. The LSTM achieves
$\text{DA} \approx 0.52$---barely above chance---while ARIMAX achieves $\text{DA} =
0.748$.

The reason is structural. Silver's predictability comes almost entirely from a
\emph{one-lag linear cross-asset signal}: yesterday's gold return, USD move, and
copper return predict today's silver return with high fidelity. This is a simple
linear relationship that ARIMAX encodes directly as a coefficient. The LSTM must
\emph{discover} the same relationship through backpropagation across noisy 20-day
sequences---with only 1,755 training observations. That is a very small dataset for a
neural network, and the LSTM does not have enough data to reliably learn a signal that
a linear model finds trivially.

The robustness checks in the notebook confirm this: the result holds across
\texttt{SEQ\_LEN} $\in \{1, 20\}$, feature sets of size 8 and 26, and learning rates
$\{10^{-3}, 10^{-4}\}$.

The key takeaway: the LSTM is not poorly implemented---it is limited by sample size
relative to the simplicity of the underlying signal.

\subsection{Sentiment in LSTM vs.\ ARIMAX}

There is one way the LSTM handles sentiment differently from ARIMAX. ARIMAX uses only
the previous day's sentiment score ($S_{t-1}$) as a single regressor. The LSTM
receives the full 20-day history of sentiment in its input sequence, so it can in
principle learn patterns like ``sustained bearish sentiment over three weeks precedes a
price drop.'' Whether it actually learns such patterns given the training set size is
an empirical question answered by comparing \texttt{LSTM-REDDIT} vs.\ \texttt{LSTM-Y}.

\subsection{Outputs saved}

\begin{itemize}
  \item \texttt{metrics\_lstm.csv} --- DA, WDA, RMSE, MAE for all four variants
  \item \texttt{preds\_lstm\_*.csv} --- actual and predicted returns for each variant
    (used for DM tests and the forecast plot)
  \item \texttt{period\_lstm\_weekly.csv} --- per-year WDA breakdown for the best
    variant (added in the period-breakdown refactor)
\end{itemize}

% ─────────────────────────────────────────────────────────────────────────────
\section{MIDAS (\texttt{weekly/03\_midas.ipynb})}
% ─────────────────────────────────────────────────────────────────────────────

\subsection{The problem MIDAS solves}

All other models in this thesis use daily data throughout. MIDAS exists to answer a
different question: can monthly macroeconomic variables (CPI, industrial production,
the federal funds rate, real rates) improve silver return forecasts, even though they
are only published once per month?

The naive approach---forward-fill the monthly value across all trading days in that
month and add it as a regressor---discards the information that there are actually
three recent monthly observations available (the last three months), each with
potentially different relevance. MIDAS handles this correctly by regressing the daily
return on a \emph{weighted sum} of multiple monthly lags, where the weights are
estimated from the data.

\subsection{The MIDAS equation}

For a single macro variable $x$ with frequency ratio $m \approx 22$ (trading days per
month) and $K = 3$ monthly lags:

\begin{equation}
  r_t = \alpha + \beta \sum_{k=1}^{K} w(k;\, \theta)\, x_{k\text{ months ago}} + \varepsilon_t
  \label{eq:midas_full}
\end{equation}

\noindent The weighting function $w(k; \theta)$ determines how much each monthly lag
contributes. The key constraint is $\sum_k w(k;\theta) = 1$, so $\beta$ captures the
total scale of the macro variable's influence while the weights capture its
\emph{decay pattern}.

\subsection{Weight functions}

\paragraph{Beta polynomial.}
\begin{equation}
  w(k;\, \theta_1, \theta_2) =
    \frac{(k/K)^{\theta_1 - 1}(1 - k/K)^{\theta_2 - 1}}{B(\theta_1, \theta_2)}
\end{equation}

The two parameters control the shape:
\begin{itemize}
  \item $\theta_1 = \theta_2 = 1$: uniform weights (same as U-MIDAS with equal
    coefficients).
  \item $\theta_1 < \theta_2$: more weight on recent lags (lag 1 dominates).
  \item Hump-shaped: lag 2 matters most if the macro signal takes a month to
    propagate into prices.
\end{itemize}

\paragraph{Exponential Almon (normalised).}
\begin{equation}
  w(k;\, \theta_1, \theta_2) = \frac{e^{\theta_1 k + \theta_2 k^2}}
    {\sum_{j=1}^{K} e^{\theta_1 j + \theta_2 j^2}}
\end{equation}

Both polynomials are estimated by minimising training MSE over the shape parameters
$(\theta_1, \theta_2)$ using the L-BFGS-B optimiser. The weight profiles are plotted
in Section~7 of the notebook, showing how far back each macro variable's influence
reaches.

\subsection{The five models}

\begin{enumerate}

\item \textbf{U-MIDAS (Unrestricted MIDAS).} Treat all $K = 3$ monthly lags of all 4
  macro variables as separate OLS regressors---$4 \times 3 = 12$ free coefficients plus
  an intercept. No constraint on the lag polynomial. This is the interpretable baseline:
  it is just OLS with many regressors.

\item \textbf{ADL-MIDAS (Autoregressive Distributed Lag).} Add yesterday's silver
  return as an AR(1) term to U-MIDAS. In practice this is the most useful variant
  because the AR term absorbs any residual autocorrelation the macro lags do not
  capture.

\item \textbf{Beta-MIDAS (single predictor).} Instead of 3 free OLS weights per
  variable, replace them with the 2-parameter Beta polynomial. Fit one model per macro
  variable. This reduces overfitting: 2 parameters instead of 3 per variable.
  The estimated $(\hat{\theta}_1, \hat{\theta}_2)$ can be interpreted directly: they
  tell you the shape of the estimated impulse response of silver returns to that macro
  variable.

\item \textbf{Multi-predictor Beta-MIDAS.} All 4 macro variables enter jointly, each
  with its own Beta polynomial. The optimiser finds 8 shape parameters simultaneously.
  This is the full MIDAS model used in the literature.

\item \textbf{MIDAS + Sentiment.} The weekly Reddit sentiment index is added as a
  fifth predictor with its own Beta polynomial. Weekly data sits at an intermediate
  frequency (between daily silver returns and monthly macro), so it enters the MIDAS
  framework naturally.

\end{enumerate}

\subsection{Why R?}

The \texttt{midasr} package provides the nonlinear optimisation for the Beta and Almon
weight functions, including gradient-based fitting of $(\theta_1, \theta_2)$ with
constraints $\theta_i > 0$. Implementing this cleanly in Python would require writing
the optimisation loop from scratch. The evaluation metrics (RMSE, MAE, DA, WDA) are
re-implemented in R at the top of the notebook to match the Python definitions exactly.

\subsection{What ``mixed frequency'' means in practice}

On any given trading day $t$, the lag matrix for CPI looks like this (with $K = 3$):

\begin{table}[H]
\centering
\caption{Lag matrix structure for CPI on a given test day}
\begin{tabular}{lll}
\toprule
Column & Contents & Recency \\
\midrule
\texttt{cpi\_lag1} & CPI published in the current month & Most recent \\
\texttt{cpi\_lag2} & CPI published in the previous month & 1 month ago \\
\texttt{cpi\_lag3} & CPI published two months ago & 2 months ago \\
\bottomrule
\end{tabular}
\end{table}

\noindent The function \texttt{build\_lag\_matrix} constructs this by, for each daily
date, looking back and grabbing the last 3 monthly observations that were available
\emph{on or before} that date. This is the correct causal construction: it uses only
information that would have been published before day $t$.

\subsection{Test window considerations}

Because monthly macro data requires at least 3 historical observations before the first
complete lag matrix can be formed, the first few rows of the test set are excluded
(this is handled by the \texttt{ok\_test} mask). In practice, with data starting in
2015 and a test set from January 2023, all test rows have complete lag matrices and the
effective test $n$ matches the other models.

\subsection{Outputs saved}

\begin{itemize}
  \item \texttt{metrics\_midas.csv} --- DA, WDA, RMSE, MAE for all model variants
  \item \texttt{preds\_umidas.csv} --- actual and predicted returns for the best
    variant (lowest RMSE)
  \item \texttt{period\_midas\_weekly.csv} --- per-year WDA breakdown (added in the
    period-breakdown refactor)
\end{itemize}

% ─────────────────────────────────────────────────────────────────────────────
\section{Daily versions}
% ─────────────────────────────────────────────────────────────────────────────

\subsection{LSTM}

The current \texttt{weekly/06\_lstm.ipynb} already uses daily data throughout
(\texttt{train.csv}/\texttt{val.csv}/\texttt{test.csv} are daily files) with
\texttt{HORIZON = 1} (predict the next \emph{day}). It is in the \texttt{weekly/}
folder by convention but is architecturally a daily model. A \texttt{daily/06\_lstm.ipynb}
would be functionally the same notebook with the output files renamed
(\texttt{metrics\_lstm\_daily.csv}, \texttt{period\_lstm\_daily.csv}) and aligned
with the daily model family.

\subsection{MIDAS}

MIDAS is already a daily-target model: the dependent variable is the daily silver
log-return $r_t$. The ``mixed'' part is the \emph{monthly} macro inputs, not the
target frequency. Creating a \texttt{daily/03\_midas.ipynb} means placing the same
model in the daily family, saving \texttt{metrics\_midas\_daily.csv} and
\texttt{period\_midas\_daily.csv}, so the evaluation notebook can include MIDAS in
the daily frequency comparison.

The notebook code does not change structurally; only the output filenames and the
\texttt{metrics\_map} entry in \texttt{evaluation.ipynb} need updating.

% ─────────────────────────────────────────────────────────────────────────────
\section{Quick reference: key differences from ARIMA/VAR/RF notebooks}
% ─────────────────────────────────────────────────────────────────────────────

\begin{table}[H]
\centering
\caption{Design comparison across model families}
\begin{tabular}{p{3cm}p{3cm}p{3cm}p{3cm}}
\toprule
 & ARIMA/VAR/RF/XGB & LSTM & MIDAS \\
\midrule
Language & Python & Python (PyTorch) & R (midasr) \\
Input frequency & Daily or weekly & Daily (20-day sequences) & Daily target, monthly inputs \\
Estimation & Closed-form or iterative & Backpropagation & NLS for weight params, OLS for scale \\
Forecast scheme & Expanding + rolling & Single pass (train/val/test split) & Single pass \\
Hyperparameter tuning & AIC / grid search & Validation loss (early stopping) & Manual ($K = 3$, Beta weights) \\
Sentiment enters as & Scalar regressor at $t-1$ & Full sequence of $S_{t-k}$ for $k = 0 \ldots 19$ & Beta-weighted MIDAS predictor \\
\bottomrule
\end{tabular}
\end{table}

\end{document}
